
\chapter{Literature Review}

The purpose of this literature review is to provide an overview of existing research on the topics of virtual reality, meditation, artificial intelligence, and physiological tools, such as biofeedback. This synthesis highlights theory and findings that form the foundation for understanding how VR can enhance mindfulness practices through emerging technology like artificial intelligence. The review is organized thematically to guide the reader from fundamental concepts such as evolution and applications of VR to areas more specific to the field of meditation and to the regions where Biofeedback mechanisms and AI-driven systems are used. These studies show how new technologies combined with meditative practices can enhance mental and emotional well-being.



 \section{Virtual Reality: Concepts and Applications}
Virtual Reality (VR) is a computer-generated environment that immerses users in realistic or imagined scenes and permits interaction through multisensory input and output \cite{jerald2016vrbook}. VR research distinguishes between immersion, describing the technical characteristics of the system, and presence, referring to the user’s subjective experience of being located within the virtual environment \cite{slater2018immersive,bryson1996vrdesign}. High levels of presence and immersion can significantly influence cognition, behavior, and emotion.

While VR initially gained prominence in gaming and entertainment, it is now widely applied in healthcare, psychotherapy, assessment, training, and education \cite{freeman2017virtual,rehm2021virtual,baker2019vrtherapy}. Meta-analytic and review work on VR in mental health indicates that VR exposure and therapeutic environments can be effective for anxiety disorders, psychosis, and trauma-related conditions when integrated into structured treatment protocols \cite{freeman2017virtual,rehm2021virtual}. Human-centered design guidelines emphasize usability, comfort, and emotional safety as key requirements for VR applications targeting vulnerable or clinical populations \cite{jerald2016vrbook}.

More recent studies have explored the integration of VR with AI and adaptive interfaces in education and training contexts, for example VR classrooms with AI-generated scenarios or intelligent tutoring and assessment \cite{kortelev2024integration,liu2024smartclassroom,elaine2025adaptivelearning}. In medical and procedural training, VR has also been used to simulate complex interventions and measure stress responses and user perceptions during high-stakes scenarios \cite{rubio2025stress}. Together, these findings position VR as a versatile medium for both skill development and emotional or stress-related interventions.

 \section{Meditation and Psychological Well-Being}
Meditation practices have been shown to influence emotional regulation and stress resilience through attentional and cognitive control mechanisms, it includes diverse practices such as mindfulness, recitation, and focused attention each offering unique physiological and cognitive benefits \cite{sedlmeier2012meta,lutz2008attention,manocha2011mind}. Evidence from the meta-analysis reveals small to moderate positive outcomes for meditation interventions on psychological stress, anxiety, depression, and overall well-being \cite{goyal2014meditation,sedlmeier2012meta}. 

Theoretical and neuroscience studies also show the potential for meditation to increase attentional power: meditation enhances attentional regulation, improve metacognitive monitoring and brain regions implicated in attentional control and emotion regulation, including the anterior cingulate cortex, insular cortex, and amygdala \cite{lutz2008attention}. Traditional conceptual models, for example, de-automatization, have explained how meditation challenges habitual processes:
Cognitive and perceptual patterns to have more adaptable and unreactive reactions to stressors \cite{deikman1966deautomatization,manocha2011mind}. 

Despite the benefits, it can be difficult for some to practice and commit to the meditation routine that is lacking due to time constraints, lack of guidance, or inability to focus attention. Consequently, recent studies have focused on technology-mediated meditation, such as smartphone apps and VR-based solutions, to offer more engaging and structured support \cite{tarrant2018virtual,cebolla2021compassion}. Evidently, these methods are aimed at filling a gap between traditional Meditation practices and modern digital life.


\section{Virtual Reality Meditation: Integrating Technology and Mindfulness}
The integration of VR and meditation, collectively referred to as VR meditation, introduces a new approach to increased relaxation and concentration in engaging, sensory-stimulating environments \cite{liszio2018nature}. In this thesis, VR meditation refers to the delivery of guided meditative practices within immersive virtual environments that combine visual, auditory, and interactive elements to support attentional focus. With the display of relaxing natural surroundings or graphical patterns along with audio, based on guidance, where it seeks to improve focus, suppress interfering distractions, and boost the subjective quality of the meditative state\cite{liszio2018nature,tarrant2018virtual}. Empirical studies tell us that immersive VR environments of nature can induce relaxation and reduce
negative affect, and improve subjective well-being compared to control conditions \cite{liszio2018nature}.

Specific protocols of VR meditation were investigated in relation to the reduction of anxiety and mood changes. improvement, where the subject undergoes guided breathing and attention training embedded in virtual nature or restorative environments \cite{tarrant2018virtual,blum2019vrbreathing}. Neuroadaptive and biofeedback-driven VR technology further refines this process by incorporating audio-visual data in real-time, as a function of physiological parameters such as Heart Rate or Respiration Rate \cite{kosunen2016neuroadaptive,weibel2020biofeedback,brief2025vrbiofeedback}. Movement-based systems such as LOOP with breath and attention to foster engagement, whereas mantra-driven meditation in VR \cite{garg2024vrmantra}.

Nonetheless, some current studies in both clinical and applied settings have expanded applications of VR meditation into various groups. For instance, VR-supported mind–body therapies have been studied in adolescents with insomnia, older persons with chronic pain, nursing students' stress levels and sleep quality  \cite{zambotti2022insomnia,sarkar2022vrpain,kim2024stress}. Virtual reality meditations, “retreats,” and supervised multi-session courses offer great opportunities; these have been suggested as scaled-up applications to provide contemplative practices to individuals who would otherwise not have attended traditional retreats \cite{ma2022vrmindfulness}. However, many of the present-day applications continue to use pre-scripted contents and will not be modified based on current physiological or emotive changes.

\section{Physiological Biofeedback and Regulation}
“Physiological biofeedback” is strongly linked with relaxation as relaxation is defined as a process where physiological responses are reduced. In this context, physiological biofeedback denotes the continuous acquisition of autonomic signals (e.g., heart rate, electrodermal activity, and peripheral temperature) to infer stress and relaxation states. The aim is to promote internal state consciousness and facilitate internal state regulation of stress and arousal. HRV biofeedback training especially focuses on breathing at a resonance frequency. This is because when a person breathes at a
frequency to best optimize the baroreflex mechanism and thus parasympathetic stimulation, causing to calmer states and better emotion regulation  \cite{lehrer2013baro}.

Electrodermal activity, usually recorded as skin conductance level (SCL) and skin conductance response (SCR), which indicates activity within the sympathetic nervous system and is a sensitive marker of emotional arousal \cite{boucsein2012eda}. In human-computer interaction and physiological computing, these include the combination of HRV, EDA, respiratory patterns, and peripheral temperature to infer mental workload, engagement, and stress  \cite{fairclough2009physio}. Incorporating these factors within the interfaces and systems enables designing closed loops to cater to users inner states.

Biofeedback systems that utilize VR can combine immersive visual or listening environments with real-time physiological feedback. There is evidence to indicate that biofeedback with virtual reality could serve to improve relaxation, presence, and self-efficacy compared to non-VR or non-biofeedback approaches \cite{weibel2020biofeedback,brief2025vrbiofeedback}. For example, virtual environments could be changed in color, brightness, or motion based on HRV or Breathing Patterns, Relaxation reinforced by intuitive metaphors \cite{chittaro2018placebo,bouny2020haptic}. Systematic reviews highlight the need for proper protocols and signal processing practices to ensure reliable results in all biofeedback studies \cite{applied2020review}.

Recent studies have also examined the application of biofeedback from emotionally challenging or stressful virtual-life experiences, for example, fear-inducing provocation or life-simulating scenarios  \cite{houzangbe2018biofeedback,rubio2025stress,rodrigues2021IMVEST}. This shows how the biological channels utilized for stress measure can be used for providing adaptive feedback support for self-regulation training.

\paragraph{Detecting Stress and Frustration through Physiological Signals.}\mbox{}\\
Stress and frustration occur when there are any cognitive, emotional, or social demands perceived as exceeding coping resources. Whereas self-report inventories, for instance, the PSS-10 and the NASA TLX, aim subjective assessments of stress and workload, physiological signals offer continuous, objective markers of autonomic nervous system activation. With acute stress, sympathetic activation will lead to acceleration of heart rate, decreased HRV, increased skin conductance, and decreased
peripheral temperature because of vasoconstriction \cite{healey2005stress,mccraty2001hrv,boucsein2012eda}.

Heart rate variability marks the constant interplay between sympathetic and parasympathetic activity; Reduced HRV is commonly interpreted as diminished parasympathetic modulation and has been associated with elevated stress and reduced self-regulatory capacity \cite{mccraty2001hrv}. 

Electrodermal activity is frequently used as a direct measure of sympathetic activation with phasic SCRs that reflect brief reactions to stimulus objects \cite{boucsein2012eda}. Skin temperature and pulse volume amplitude (PVA) offer complementary information regarding peripheral blood flow and vasomotor tone, which also varies depending on stress and frustration \cite{setz2010frustration}.

In human-computer interaction and gaming studies, these multimodal variables are used on
to distinguish between stress and cognitive load and to evaluate the effect of interactive tasks on user \cite{fairclough2009physio,aliyari2018gamestress,porter2019GamesStress}. Standardized laboratory procedures such as the Trier Social Stress Test (TSST) and its VR adaptations illustrate how psychophysiological responses to stress can be reliably elicited and measured under controlled conditions \cite{Kirschbaum1993TSST,Allen2014TSSTReview,Zimmer2019VR_TSST,OpenTSSTVR2020,rodrigues2021IMVEST}. 

Devices such as the \textbf{Biofeedback 2000 Xpert} enable synchronized collection of multiple biosignals: BVP-derived heart rate, SCL, skin temperature, and pulse frequency, making them suitable for both stress-induction tasks and VR-based relaxation training. Combining continuous physiological data with behavioral measures and questionnaires offers a comprehensive view of stress dynamics during tasks such as game-based stress induction and subsequent meditation.



\section{Artificial Intelligence in Adaptive Virtual Environments}
Artificial intelligence (AI) technology is transforming world of virtual reality (VR) by introducing adaptivity, personalization, and intelligent guidance \cite{zhou2020adaptive}. In VR, AI can be used to model user skills, detect affective states, or personalize narrative and feedback. Adaptive interfaces and agent-based systems have been proposed to tailor VR experiences to individual needs, improving usability, learning, and engagement \cite{zhou2020adaptive,arntz2019aivr}.

For example, In therapeutic and training applications, AI–VR systems have been used to personalize social and cognitive training, including support for autism and other neuro-developmental conditions \cite{bhatt2020autism}. Conversational agents and AI-driven characters in VR may model their social interactions, offer guidance, and adjust their responses based on user behaviour or physiological signals \cite{schlagowski2021empathy}. Multimodal AI systems that interpret gaze, posture, and physiological markers can further refine personalization and adaptivity \cite{martin2019multimodal}.

Recent developments in AI-enhanced VR has focused on immersive training and simulation for real-world skill development, underscoring the potential of AI to modify scenario difficulty, feedback, and pacing in real time \cite{elaine2025adaptivelearning}. Other studies discuss challenges in quantifying immersion and workload in complex VR tasks and highlight the role of secondary tasks and physiological measures in evaluating adaptive systems \cite{payne2024immersion,rubio2025stress}. These developments motivate the integration of AI-driven guidance and decision rules into VR meditation systems that can adapt prompts and sensory parameters to users’ momentary stress level or self-reported state.

\section{Multimodal Interaction and System Design}
Immersive experiences are most effective when they engage multiple senses simultaneously. Multimodal VR systems combine audio, visual, tactile, and sometimes olfactory feedback to enhance emotional presence and relaxation depth \cite{aimagery2023olfactory,martin2019multimodal}.  In the context of meditation and relaxation, soft lighting, gentle motion, ambient music, and nature sounds can be orchestrated to support calm, focused attention \cite{blum2019vrbreathing,shaw2011meditationchamber}.

Research on multimodal interaction in VR emphasizes that carefully balancing sensory input is crucial: overly complex or intense stimuli can increase cognitive load and undermine relaxation, whereas minimalistic but meaningful cues can facilitate mindfulness and comfort \cite{martin2019multimodal,freeman2017virtual}. Recent work explores multisensory anxiety-reduction approaches that incorporate AI, olfactory stimuli, and biofeedback-enhanced guided imagery, demonstrating the potential of combining smell, sound, and visual imagery in immersive anxiety-reduction protocols \cite{aimagery2023olfactory}. Similarly, biofeedback game mechanics in VR, where environmental parameters change based on fear or arousal, highlight how physiological signals can be embedded into the interaction loop \cite{houzangbe2018biofeedback}.

Studies on immersion measurement suggest that secondary tasks and performance metrics can be used alongside self-report presence scales to quantify how deeply users are engaged in multimodal VR environments \cite{payne2024immersion}. Effective multimodal design for VR meditation thus lies at the intersection of aesthetics, psychophysiology, and adaptive engineering.

\section{Psychometric Measures Used in the Study}

% To evaluate the effectiveness of VR meditation and biofeedback interventions, several validated psychometric instruments were selected. These measures assess baseline stress, mindfulness, affect, and post-session experience across cognitive, emotional, and perceptual domains.
To complement physiological measurements, this study employs validated psychometric instruments to assess perceived stress, mindfulness, affect, presence, workload, cybersickness, and user experience.

The \textbf{Perceived Stress Scale (PSS-10)} captures the degree to which individuals appraise situations in their lives as unpredictable, uncontrollable, and overloading \cite{cohen1983pss}. The \textbf{Mindful Attention Awareness Scale (MAAS)} measures dispositional mindfulness, defined as open or receptive awareness of and attention to present-moment experience \cite{brown2003maas}. 

Emotional valence is assessed using the \textbf{Positive and Negative Affect Schedule (PANAS)}, which includes two independent subscales for positive and negative affect \cite{watson1988panas}. State mindfulness during or immediately after the meditation session is measured with the \textbf{State Mindfulness Scale (SMS)}, which captures awareness of mind and body \cite{tanay2013sms}. 

In the VR condition, presence is evaluated with the \textbf{I-Group Presence Questionnaire (IPQ)}, which measures spatial presence, involvement, and experienced realism \cite{schubert2001ipq}. Cybersickness symptoms are assessed using the \textbf{CSQ-VR}, a validated questionnaire covering nausea, disorientation, and oculomotor discomfort in VR environments \cite{kourtesis2019csqvr}. 

Perceived workload is measured with the \textbf{NASA Task Load Index (NASA-TLX)}, providing ratings along six dimensions: mental, physical, and temporal demand, performance, effort, and frustration \cite{hart1988nasatlx}. Overall usability and experiential quality are evaluated with the \textbf{User Experience Questionnaire—Short (UEQ-S)}, which yields pragmatic and hedonic quality scores \cite{schrepp2017ueqs}. Together, these instruments provide a multidimensional assessment of stress reduction, mindfulness induction, emotional state, cognitive workload, and experiential quality.

\section{Commercial Applications and Research Gap}
Commercial VR meditation platforms such as TRIPP, Guided Meditation VR, and FlowVR have made immersive mindfulness tools more publicly accessible and illustrate the growing interest in VR for emotional self-care and stress management \cite{ghosal2025vrmindfulness}. Pilot trials with youth and other vulnerable groups show that VR meditation can be feasible and acceptable, and may reduce anxiety or distress in challenging life situations \cite{chavez2020youth}. Clinical and applied studies further indicate that VR-based meditative interventions can help modulate arousal and improve sleep or pain-related outcomes \cite{zambotti2022insomnia,sarkar2022vrpain,kim2024stress}.

However, reviews of commercial and research prototypes alike note that most VR meditation systems still rely on static scenarios and linear, pre-scripted guidance \cite{ghosal2025vrmindfulness}. They rarely integrate real-time physiological monitoring or AI-driven personalisation, and they often lack mechanisms for adapting content to users’ momentary stress level, workload, or frustration. Recent work on VR stress measurement and training in medical and educational contexts emphasizes the importance of combining psychophysiological measures with user-centered design to truly understand and support users under stress \cite{rubio2025stress,rodrigues2021IMVEST}.

This literature therefore identifies a clear research gap: the need for AI-guided, biofeedback-informed VR meditation systems that can adapt instructions and audiovisual parameters dynamically based on physiological and self-report indicators. Such systems have the potential to deliver more individualized, evidence-based relaxation experiences and to advance the clinical reliability and scalability of digital mindfulness interventions.

% \section{Summary} 
In summary, VR provides immersive, controllable environments capable of supporting mental health interventions, including meditation and stress management. Meditation itself has a robust evidence base for improving psychological well-being, and VR-based meditation systems extend these practices into interactive, engaging formats. Physiological biofeedback, particularly through HRV and electrodermal activity, offers objective insight into stress and relaxation, while AI and multimodal design enable adaptive, personalized VR experiences.

Despite promising developments in VR, current VR meditation applications rarely combine AI-driven guidance with real-time physiological adaptation. Addressing this gap, the present work focuses on the design and evaluation of an AI-guided, biofeedback-based VR meditation system that aims to reduce stress and enhance emotional regulation by responding dynamically to users’ physiological and self-reported states.









